% ============================================================
%  基于AHP-FCE模型的公司综合实力评价研究
%  数学建模论文 · XeLaTeX 编译
% ============================================================
\documentclass[a4paper,12pt]{ctexart}

% ---------- 页面布局 ----------
\usepackage[top=2.5cm, bottom=2.5cm, left=2.8cm, right=2.8cm]{geometry}
\usepackage{setspace}
\setstretch{1.3}

% ---------- 数学与公式 ----------
\usepackage{amsmath,amssymb}

% ---------- 图表 ----------
\usepackage{graphicx}
\usepackage{float}
\usepackage{booktabs}
\usepackage{array}
\usepackage{multirow}
\usepackage{subcaption}
\usepackage{caption}
\captionsetup[figure]{font=small,labelfont=bf,labelsep=quad}
\captionsetup[table]{font=small,labelfont=bf,labelsep=quad}

% ---------- 交叉引用 ----------
\usepackage{hyperref}
\hypersetup{colorlinks=true,linkcolor=blue,citecolor=blue,urlcolor=blue}
\usepackage{cleveref}
\crefname{equation}{式}{式}
\crefname{figure}{图}{图}
\crefname{table}{表}{表}
\crefname{section}{节}{节}

% ---------- 页眉页脚 ----------
\usepackage{fancyhdr}
\pagestyle{fancy}
\fancyhf{}
\fancyhead[C]{\kaishu\small 基于AHP-FCE模型的公司综合实力评价研究}
\fancyfoot[C]{\thepage}
\renewcommand{\headrulewidth}{0.4pt}

% ---------- 参考文献 ----------
\usepackage{cite}

% ---------- 列表 ----------
\usepackage{enumitem}

\begin{document}

% ============================================================
%  标题页
% ============================================================
\title{{\heiti\zihao{2} 基于AHP-FCE模型的公司综合实力评价研究}\\[6pt]
       {\kaishu\zihao{-3} ——模糊综合评价与层次分析法的多指标决策应用}}
\author{{\kaishu\zihao{4} 冯浩然}}
\date{}
\maketitle

\thispagestyle{empty}

% ============================================================
%  摘要
% ============================================================
\begin{abstract}
\noindent
针对公司综合实力评价中指标多源异构、边界模糊及主观判断难以量化的问题，
本文构建了基于层次分析法（AHP）与模糊综合评价（FCE）相结合的多指标决策模型。
首先，建立了包含产品综合评价、成本费用和销售能力3个一级指标及19个二级指标的三层评价体系，
并据此构造了七家公司的模拟样本数据；
其次，设定五级评价等级$V=\{\text{优秀},\text{良好},\text{中等},\text{较差},\text{很差}\}$及其分值向量
$S=[100,80,60,40,20]$，采用梯形模糊分布建立隶属函数实现指标值到等级隶属度的模糊转化，
利用AHP方根法确定各级权重并通过一致性检验（全部$CR<0.1$）；
最后，基于$M(\cdot,\oplus)$加权平均算子对每家公司独立执行两级模糊合成，
将得到的五等级隶属度向量与等级分值加权求和，转化为综合得分并排序。
实证结果显示：公司d综合得分最高（81.14分，等级"优秀"），其后依次为
a（78.15）、c（71.02）、e（70.50）、b（68.66）、f（64.79）、g（63.04）。
该方法有效保留了指标值在等级边界处的模糊过渡信息，评价结果合理、可解释，
为企业竞争力评估与标杆管理提供了系统的定量分析框架。

\vspace{0.8em}
\noindent{\textbf{关键词：}} 模糊综合评价；层次分析法；梯形隶属函数；多指标决策；公司综合实力评价
\end{abstract}

% ============================================================
%  1 引言
% ============================================================
\section{引言}

\subsection{问题背景与研究意义}

企业综合实力的科学评价是投资者决策、供应链选优与政府监管的关键依据。
然而，企业评价面临三重挑战：
其一，\textbf{指标多源异构}——评价维度涵盖产品品质、成本控制、销售能力等多个方面，
指标量纲与方向各异；
其二，\textbf{等级边界模糊}——"优秀"与"良好"之间并无明确的数值界线，
传统评分法将评价视为刚性分类，忽略了过渡信息；
其三，\textbf{权重主观性强}——各项指标对综合实力的贡献不同，
如何将专家经验转化为逻辑一致的量化权重是建模难点。

\subsection{研究思路}

针对上述问题，本文采用模糊综合评价（Fuzzy Comprehensive Evaluation, FCE）
与层次分析法（Analytic Hierarchy Process, AHP）相结合的建模策略。
FCE通过隶属函数将精确指标值转化为对各评价等级的隶属度，
保留了指标在等级过渡带的全部信息\cite{wang1983}；
AHP通过两两比较矩阵将主观判断转化为具有一致性保证的数值权重\cite{saaty1980}。
二者结合兼具"处理模糊性"与"系统化赋权"的优势。

本文以七家公司（a$\sim$g）为评价对象，按以下主线展开：
（1）构建三层评价指标体系；（2）设计模拟样本数据并完成标准化；
（3）设定五级评价标尺，建立梯形隶属函数和AHP多层权重体系；
（4）对每家公司独立执行两级模糊合成，计算综合得分并排序；
（5）通过多维可视化图表展示评价结果。

\section{评价指标体系构建}

\subsection{设计原则}

指标体系的构建遵循四项原则：
\textbf{系统性}——覆盖企业运营的主要方面；
\textbf{可操作性}——指标数据可获取、可量化；
\textbf{层次性}——指标按层级组织以反映不同粒度的评价信息；
\textbf{独立性}——同层指标间信息重叠尽可能小。

\subsection{三层体系结构}

本文构建"目标层—准则层—指标层"三层评价体系。
目标层（A）为公司综合实力，是评价的最终对象。
准则层（B）包含3个一级指标，指标层（C）包含19个二级指标，
体系的树形结构如\cref{fig:structure}所示（见下页）。

\begin{figure}[H]
\centering
\begin{minipage}{0.82\textwidth}
\centering
{\heiti\large 目标层（A）：公司综合实力}

\vspace{0.4cm}
\begin{tabular}{c}
\hline
{\heiti 公司综合实力评价 A} \\[2pt]
\hline
\end{tabular}

\vspace{0.7cm}
\begin{minipage}[t]{0.30\textwidth}
\centering
\begin{tabular}{|c|}
\hline
{\heiti B$_1$ 产品综合评价} \\[2pt]
\hline
C$_1$ 产品品牌 \\
C$_2$ 产品质量 \\
C$_3$ 性价比 \\
C$_4$ 产品款式 \\
C$_5$ 产品包装 \\
C$_6$ 店铺信用度 \\
C$_7$ 售后服务 \\
\hline
\end{tabular}
\end{minipage}
\hfill
\begin{minipage}[t]{0.30\textwidth}
\centering
\begin{tabular}{|c|}
\hline
{\heiti B$_2$ 成本费用} \\[2pt]
\hline
C$_8$ 材料费 \\
C$_9$ 运输费用 \\
C$_{10}$ 设备维修费用 \\
C$_{11}$ 运行费用 \\
C$_{12}$ 设备折旧费用 \\
C$_{13}$ 人员工资 \\
C$_{14}$ 电耗费用 \\
C$_{15}$ 水耗费用 \\
\hline
\end{tabular}
\end{minipage}
\hfill
\begin{minipage}[t]{0.30\textwidth}
\centering
\begin{tabular}{|c|}
\hline
{\heiti B$_3$ 销售能力} \\[2pt]
\hline
C$_{16}$ 销售人员能力 \\
C$_{17}$ 广告宣传 \\
C$_{18}$ 行业需求 \\
C$_{19}$ 家庭需求 \\
\hline
\end{tabular}
\end{minipage}

\vspace{0.4cm}
{\kaishu 准则层（B）：3个一级指标}

\vspace{0.2cm}
{\kaishu 指标层（C）：19个二级指标}
\end{minipage}
\caption{三层评价指标体系结构}
\label{fig:structure}
\end{figure}

\subsection{指标分类}

19个二级指标按其数值方向分为三类：

\begin{enumerate}[label=(\arabic*),itemsep=2pt]
    \item \textbf{偏大型（效益型，9项）：}数值越大评价越优。
          包括C$_2$产品质量、C$_3$性价比、C$_4$产品款式、C$_6$店铺信用度、
          C$_7$售后服务、C$_{16}$销售人员能力、C$_{17}$广告宣传、
          C$_{18}$行业需求、C$_{19}$家庭需求。
    \item \textbf{偏小型（成本型，8项）：}数值越小评价越优。
          包括C$_8$材料费、C$_9$运输费用、C$_{10}$设备维修费用、
          C$_{11}$运行费用、C$_{12}$设备折旧费用、C$_{13}$人员工资、
          C$_{14}$电耗费用、C$_{15}$水耗费用。
    \item \textbf{中间型（适度型，2项）：}数值落在最优区间内评价最优。
          包括C$_1$产品品牌、C$_5$产品包装。
\end{enumerate}

\section{数据来源与预处理}

\subsection{数据构造}

本文以七家公司（a$\sim$g）为评价对象。由于实际企业数据涉及商业保密且方法演示需要
覆盖不同的公司特征，本文采用\textbf{模拟样本数据}。
数据的构造遵循三项准则：
（1）各指标取值符合实际经济含义的取值范围；
（2）公司间具有差异化特征，避免评价结果高度同质化；
（3）同一公司内不同指标间逻辑自洽（如成本优势型企业的各项成本指标应一致偏低）。

全部数据存储于 \texttt{company\_data.xlsx}，含"原始数据""数据（仅数值）""指标体系"三个工作表，
代码实现基于 Python（NumPy + Pandas + Matplotlib）。

\subsection{评分标准化}

为消除量纲差异、统一评价标尺，将19个指标映射到0$\sim$100分制：

\begin{enumerate}[label=(\arabic*),itemsep=2pt]
    \item \textbf{偏大型指标：}以行业通用的百分制评分。如产品质量反映该企业在同类市场中的
          品质排名（100分为业界标杆），性价比衡量功能与价格的匹配程度。取值区间约60$\sim$100。
    \item \textbf{偏小型指标：}以成本费用占营业收入的百分比经线性映射得到。
          如材料费得分 =（材料费$/$营业收入）$\times$ 100。
          成本占比越低则得分越低（即越优）。取值区间约5$\sim$40。
    \item \textbf{中间型指标：}以与行业最优区间的偏离程度评分，接近最优区间中心者得高分。
          取值区间约65$\sim$89。
\end{enumerate}

\subsection{数据概览}

标准化后的数据矩阵（7公司 $\times$ 19指标）如\cref{tab:rawdata}所示，
该表是后续所有计算的唯一输入。

\begin{table}[H]
\centering
\caption{七家公司19项评价指标标准化数据（0$\sim$100分制）}
\label{tab:rawdata}
\scriptsize
\begin{tabular}{lccccccc}
\toprule
指标 & a & b & c & d & e & f & g \\
\midrule
C$_1$ 产品品牌      & 81.0 & 82.3 & 66.7 & 89.8 & 77.8 & 88.6 & 88.4 \\
C$_2$ 产品质量      & 94.7 & 88.6 & 77.8 & 100.0& 87.9 & 79.4 & 83.2 \\
C$_3$ 性价比        & 87.0 & 76.6 & 73.0 & 84.4 & 85.3 & 75.6 & 70.7 \\
C$_4$ 产品款式      & 89.1 & 71.7 & 74.4 & 81.2 & 75.1 & 62.3 & 68.8 \\
C$_5$ 产品包装      & 65.3 & 68.4 & 77.2 & 80.8 & 69.3 & 73.1 & 76.7 \\
C$_6$ 店铺信用度    & 100.0& 76.5 & 84.6 & 92.6 & 81.6 & 72.1 & 77.8 \\
C$_7$ 售后服务      & 86.9 & 78.8 & 81.1 & 93.2 & 78.1 & 73.7 & 70.9 \\
C$_8$ 材料费        & 21.8 & 21.8 & 27.3 & 8.6  & 38.4 & 37.9 & 38.1 \\
C$_9$ 运输费用      & 13.2 & 15.1 & 18.3 & 6.2  & 24.5 & 27.8 & 27.7 \\
C$_{10}$ 设备维修费用& 8.6  & 5.1  & 12.6 & 6.5  & 16.7 & 19.0 & 22.0 \\
C$_{11}$ 运行费用    & 7.8  & 18.8 & 16.1 & 13.2 & 23.8 & 26.7 & 28.2 \\
C$_{12}$ 设备折旧费用& 15.7 & 8.1  & 19.8 & 6.2  & 19.0 & 23.9 & 17.5 \\
C$_{13}$ 人员工资    & 25.8 & 27.5 & 26.8 & 18.3 & 33.9 & 30.1 & 40.0 \\
C$_{14}$ 电耗费用    & 6.5  & 4.8  & 11.6 & 4.0  & 15.1 & 17.7 & 13.8 \\
C$_{15}$ 水耗费用    & 5.1  & 6.3  & 9.4  & 0.6  & 10.7 & 9.1  & 8.2  \\
C$_{16}$ 销售人员能力& 78.0 & 68.6 & 83.2 & 88.3 & 65.2 & 76.6 & 73.9 \\
C$_{17}$ 广告宣传    & 73.2 & 61.0 & 76.6 & 85.7 & 61.8 & 70.2 & 58.0 \\
C$_{18}$ 行业需求    & 70.1 & 74.4 & 90.1 & 79.3 & 71.8 & 76.1 & 65.7 \\
C$_{19}$ 家庭需求    & 74.4 & 67.0 & 86.0 & 82.5 & 63.3 & 79.5 & 68.6 \\
\bottomrule
\end{tabular}
\end{table}

从\cref{tab:rawdata}可观察各公司的差异化特征：
公司d的成本费用指标（C$_8$$\sim$C$_{15}$）普遍低于其他公司，
尤其是材料费（8.6）和水耗费用（0.6）远低于行业均值，是典型的成本优势型企业；
公司a的产品维度指标（C$_1$$\sim$C$_7$）均衡且高分，多数在90分以上；
公司在成本维度的明显分化为后续评价提供了充分的区分度。

\section{模糊综合评价模型}

在完成指标体系构建与数据准备后，本章依次建立模型的三个核心组件：
评价等级与隶属函数（将精确值模糊化）、AHP权重体系（确定各指标重要性）、
两级模糊合成流程（组合隶属度与权重得到最终评价）。

\subsection{评价等级设定}

模糊综合评价需要预先设定\textbf{评价标尺}。
本文参考工程模糊集理论的通用框架\cite{chen2000}，将评价等级设为五级：

\begin{equation}
    V = \{\text{优秀}(v_1),\ \text{良好}(v_2),\ \text{中等}(v_3),\
           \text{较差}(v_4),\ \text{很差}(v_5)\}
\end{equation}

每个等级分配一个代表分值，采用等距递减方式与百分制对应：

\begin{equation}
    S = [s_1,\ s_2,\ s_3,\ s_4,\ s_5] = [100,\ 80,\ 60,\ 40,\ 20]
\end{equation}

选择五等级的理由有三：
（1）与百分制下每20分一档自然对应；
（2）在区分度与评判者认知负荷间取得平衡——三等级过于粗糙，七等级以上相邻等级语义重叠严重；
（3）五等级是FCE领域最成熟的设级方式，大量隶属函数参数研究均基于此框架\cite{xie2000}。

\subsection{梯形隶属函数}

隶属函数是连接"精确指标值"与"模糊评价等级"的桥梁。
本文采用\textbf{梯形模糊分布}，其由四个参数$[a,b,c,d]$定义：
在$[b,c]$区间内隶属度为1（完全隶属），在$[a,b]$和$[c,d]$区间内线性过渡，
区间外为0。梯形分布兼具"稳定隶属平台"与"平滑过渡斜坡"，
几何意义直观、计算简便。

\subsubsection{三种函数形式}

根据指标类型，使用三种方向不同的梯形函数。

\textbf{偏大型（升半梯形）：}适用于效益型指标，指标值越大隶属度越高。

\begin{equation}
\mu_{\text{up}}(x; a, b, c, d) =
\begin{cases}
0, & x \leq a \\[4pt]
\dfrac{x - a}{b - a}, & a < x \leq b \\[8pt]
1, & b < x \leq c \\[4pt]
\dfrac{d - x}{d - c}, & c < x \leq d \\[8pt]
0, & x > d
\end{cases}
\label{eq:up}
\end{equation}

\textbf{偏小型（降半梯形）：}适用于成本型指标，指标值越小隶属度越高。

\begin{equation}
\mu_{\text{down}}(x; a, b, c, d) =
\begin{cases}
1, & x \leq a \\[4pt]
\dfrac{b - x}{b - a}, & a < x \leq b \\[8pt]
0, & b < x \leq c \\[4pt]
\dfrac{x - c}{d - c}, & c < x \leq d \\[8pt]
0, & x > d
\end{cases}
\label{eq:down}
\end{equation}

\textbf{中间型（全梯形）：}适用于适度型指标，在中间区间隶属度最高。

\begin{equation}
\mu_{\text{mid}}(x; a, b, c, d) =
\begin{cases}
0, & x \leq a \\[4pt]
\dfrac{x - a}{b - a}, & a < x \leq b \\[8pt]
1, & b < x \leq c \\[4pt]
\dfrac{d - x}{d - c}, & c < x \leq d \\[8pt]
0, & x > d
\end{cases}
\label{eq:mid}
\end{equation}

\subsubsection{参数设定}

将五级评价标尺与三种梯形函数结合，需为每个"等级$\times$指标类型"组合设定一组$[a,b,c,d]$参数。
参数设置的核心原则是\textbf{相邻等级之间具有合理的过渡重叠区}，
确保每个指标值对至少一个等级具有非零隶属度，同时保持等级区分度。

以偏大型指标的"优秀"与"良好"为例：
优秀参数$[80,90,100,100]$，良好参数$[70,80,80,90]$。
两者在80$\sim$90区间形成对称过渡——在85分处，对"优秀"的隶属度为$(85-80)/10 = 0.5$，
对"良好"的隶属度为$(90-85)/10 = 0.5$，恰好各半，直观体现了"边界模糊"特性。
全部15组参数（5等级$\times$3类型）如\cref{tab:params}所示。

\begin{table}[H]
\centering
\caption{五等级隶属函数参数 $[a,b,c,d]$（5等级 $\times$ 3指标类型 = 15组）}
\label{tab:params}
\begin{tabular}{cccccc}
\toprule
\multirow{2}{*}{评价等级} & \multirow{2}{*}{分值} &
\multicolumn{3}{c}{指标类型} \\
\cmidrule(lr){3-5}
& & 偏大型 & 偏小型 & 中间型 \\
\midrule
$v_1$（优秀） & 100 & (80, 90, 100, 100) & (0, 0, 15, 25)  & (65, 75, 85, 95) \\
$v_2$（良好） & 80  & (70, 80, 80, 90)   & (15, 25, 25, 35) & (55, 65, 65, 75) \\
$v_3$（中等） & 60  & (60, 70, 70, 80)   & (25, 35, 35, 45) & (45, 55, 55, 65) \\
$v_4$（较差） & 40  & (50, 60, 60, 70)   & (35, 45, 45, 55) & (35, 45, 45, 55) \\
$v_5$（很差） & 20  & (0, 0, 50, 60)     & (45, 55, 100, 100) & (0, 0, 35, 45) \\
\bottomrule
\end{tabular}
\end{table}

\subsection{层次分析法（AHP）确定权重}

\subsubsection{方法原理}

AHP由Saaty于1977年提出，核心是将复杂的多因素权重问题转化为一系列两两比较判断，
再通过数学方法从判断矩阵中提取权重向量\cite{saaty1980}。
本文采用方根法（几何平均法），其计算量低于特征向量法而精度几乎相当。

\textbf{步骤1——构造判断矩阵：}采用1$\sim$9标度法（\cref{tab:scale}），
对同层指标两两比较，构建正互反矩阵$A=(a_{ij})_{n\times n}$，
满足$a_{ij}>0$，$a_{ji}=1/a_{ij}$，$a_{ii}=1$。

\begin{table}[H]
\centering
\caption{1$\sim$9标度含义}
\label{tab:scale}
\begin{tabular}{cl}
\toprule
标度 & 含义 \\
\midrule
1   & 两个因素同等重要 \\
3   & 前者比后者稍微重要 \\
5   & 前者比后者明显重要 \\
7   & 前者比后者强烈重要 \\
9   & 前者比后者极端重要 \\
2,4,6,8 & 相邻判断的中间值 \\
倒数   & 因素$j$与$i$比较 $a_{ji}=1/a_{ij}$ \\
\bottomrule
\end{tabular}
\end{table}

\textbf{步骤2——方根法求权重：}

\begin{equation}
    w_i = \frac{\bigl(\prod_{j=1}^{n} a_{ij}\bigr)^{1/n}}
                {\sum_{i=1}^{n}\bigl(\prod_{j=1}^{n} a_{ij}\bigr)^{1/n}},\quad i=1,\ldots,n
\label{eq:geometric}
\end{equation}

\textbf{步骤3——一致性检验：}

\begin{equation}
    \lambda_{\max} = \frac{1}{n}\sum_{i=1}^{n}\frac{(AW)_i}{w_i},\quad
    CI = \frac{\lambda_{\max} - n}{n - 1},\quad
    CR = \frac{CI}{RI}
\label{eq:cr}
\end{equation}

其中$RI$为随机一致性指标（\cref{tab:ri}）。当$CR < 0.1$时通过检验，
表明判断矩阵具有满意的一致性，权重结果逻辑可靠。

\begin{table}[H]
\centering
\caption{随机一致性指标$RI$}
\label{tab:ri}
\begin{tabular}{cccccccccc}
\toprule
$n$ & 1 & 2 & 3 & 4 & 5 & 6 & 7 & 8 & 9 & 10 \\
\midrule
$RI$ & 0.00 & 0.00 & 0.58 & 0.90 & 1.12 & 1.24 & 1.32 & 1.41 & 1.45 & 1.49 \\
\bottomrule
\end{tabular}
\end{table}

\subsubsection{一级指标权重（A-B层）}

产品综合评价（B$_1$）是企业的核心竞争力，重要程度高于成本费用（B$_2$）；
销售能力（B$_3$）受产品与成本的双重制约，在三者中权重最低。
据此构造判断矩阵如\cref{tab:ab}，权重为
$W_B = [0.5396,\ 0.2970,\ 0.1634]$，
$CR = 0.0079 < 0.1$，一致性良好。

\begin{table}[H]
\centering
\caption{一级指标判断矩阵及权重}
\label{tab:ab}
\begin{tabular}{c|ccc|c}
\toprule
A & B$_1$（产品） & B$_2$（成本） & B$_3$（销售） & 权重 \\
\midrule
B$_1$ 产品综合评价 & 1   & 2   & 3   & 0.5396 \\
B$_2$ 成本费用     & 1/2 & 1   & 2   & 0.2970 \\
B$_3$ 销售能力     & 1/3 & 1/2 & 1   & 0.1634 \\
\bottomrule
\multicolumn{5}{l}{$\lambda_{\max}=3.0092$，$CI=0.0046$，$CR=0.0079<0.1$} \\
\end{tabular}
\end{table}

\subsubsection{二级指标权重（B-C层）}

\textbf{B$_1$产品综合评价下的7个指标（\cref{tab:b1c}）：}
产品质量（C$_2$）和性价比（C$_3$）是消费决策的核心依据，权重最高且相等；
店铺信用度（C$_6$）和售后服务（C$_7$）影响购买信心和复购率，权重次之；
品牌（C$_1$）、款式（C$_4$）、包装（C$_5$）属附加价值，权重最低。

\begin{table}[H]
\centering
\caption{B$_1$产品综合评价下属二级指标判断矩阵}
\label{tab:b1c}
\small
\begin{tabular}{c|ccccccc|c}
\toprule
B$_1$ & C$_1$ & C$_2$ & C$_3$ & C$_4$ & C$_5$ & C$_6$ & C$_7$ & 权重 \\
\midrule
C$_1$ 品牌   & 1 & 1/3 & 1/3 & 1   & 1   & 1/2 & 1/2 & 0.0758 \\
C$_2$ 质量   & 3 & 1   & 1   & 3   & 3   & 2   & 2   & 0.2468 \\
C$_3$ 性价比 & 3 & 1   & 1   & 3   & 3   & 2   & 2   & 0.2468 \\
C$_4$ 款式   & 1 & 1/3 & 1/3 & 1   & 1   & 1/2 & 1/2 & 0.0758 \\
C$_5$ 包装   & 1 & 1/3 & 1/3 & 1   & 1   & 1/2 & 1/2 & 0.0758 \\
C$_6$ 信用   & 2 & 1/2 & 1/2 & 2   & 2   & 1   & 1   & 0.1396 \\
C$_7$ 售后   & 2 & 1/2 & 1/2 & 2   & 2   & 1   & 1   & 0.1396 \\
\bottomrule
\multicolumn{9}{l}{$\lambda_{\max}=7.0203$，$CI=0.0034$，$CR=0.0026<0.1$} \\
\end{tabular}
\end{table}

\textbf{B$_2$成本费用下的8个指标（\cref{tab:b2c}）：}
材料费（C$_8$）是制造业最大成本项，权重最高；
运输、维修、运行、折旧（C$_9$$\sim$C$_{12}$）属运营支出，权重次之且相等；
人员工资（C$_{13}$）再次；电耗与水耗（C$_{14}$、C$_{15}$）占比最小。

\begin{table}[H]
\centering
\caption{B$_2$成本费用下属二级指标判断矩阵}
\label{tab:b2c}
\small
\begin{tabular}{c|cccccccc|c}
\toprule
B$_2$ & C$_8$ & C$_9$ & C$_{10}$ & C$_{11}$ & C$_{12}$ & C$_{13}$ & C$_{14}$ & C$_{15}$ & 权重 \\
\midrule
C$_8$ 材料费       & 1   & 2   & 2   & 2   & 2   & 3   & 4   & 4   & 0.2499 \\
C$_9$ 运输费       & 1/2 & 1   & 1   & 1   & 1   & 2   & 3   & 3   & 0.1433 \\
C$_{10}$ 维修费    & 1/2 & 1   & 1   & 1   & 1   & 2   & 3   & 3   & 0.1433 \\
C$_{11}$ 运行费    & 1/2 & 1   & 1   & 1   & 1   & 2   & 3   & 3   & 0.1433 \\
C$_{12}$ 折旧费    & 1/2 & 1   & 1   & 1   & 1   & 2   & 3   & 3   & 0.1433 \\
C$_{13}$ 工资      & 1/3 & 1/2 & 1/2 & 1/2 & 1/2 & 1   & 2   & 2   & 0.0798 \\
C$_{14}$ 电耗      & 1/4 & 1/3 & 1/3 & 1/3 & 1/3 & 1/2 & 1   & 1   & 0.0485 \\
C$_{15}$ 水耗      & 1/4 & 1/3 & 1/3 & 1/3 & 1/3 & 1/2 & 1   & 1   & 0.0485 \\
\bottomrule
\multicolumn{10}{l}{$\lambda_{\max}=8.0413$，$CI=0.0059$，$CR=0.0042<0.1$} \\
\end{tabular}
\end{table}

\textbf{B$_3$销售能力下的4个指标（\cref{tab:b3c}）：}
销售人员能力（C$_{16}$）直接决定销售执行效率，权重最高；
行业需求（C$_{18}$）与家庭需求（C$_{19}$）反映市场容量，两者等权；
广告宣传（C$_{17}$）为辅助手段，权重最低。

\begin{table}[H]
\centering
\caption{B$_3$销售能力下属二级指标判断矩阵}
\label{tab:b3c}
\begin{tabular}{c|cccc|c}
\toprule
B$_3$ & C$_{16}$ & C$_{17}$ & C$_{18}$ & C$_{19}$ & 权重 \\
\midrule
C$_{16}$ 销售人员能力 & 1   & 3   & 2   & 2   & 0.4231 \\
C$_{17}$ 广告宣传     & 1/3 & 1   & 1/2 & 1/2 & 0.1222 \\
C$_{18}$ 行业需求     & 1/2 & 2   & 1   & 1   & 0.2274 \\
C$_{19}$ 家庭需求     & 1/2 & 2   & 1   & 1   & 0.2274 \\
\bottomrule
\multicolumn{6}{l}{$\lambda_{\max}=4.0104$，$CI=0.0035$，$CR=0.0038<0.1$} \\
\end{tabular}
\end{table}

四个矩阵的$CR$均在$10^{-3}$$\sim$$10^{-2}$量级，远低于0.1阈值，权重体系具有高度逻辑一致性。

\subsubsection{综合权重合成}

二级指标对总目标的综合权重由$w_i^{\text{综合}} = w_B^{(i)} \cdot w_C^{(i)}$计算后经全体归一化得到。
全部19个指标的综合权重及其排序如\cref{fig:weight}所示（完整数据见\cref{tab:weights}）。

\begin{figure}[H]
\centering
\includegraphics[width=0.92\textwidth]{weight_distribution.png}
\caption{19个评价指标综合权重分布（AHP方根法 + 两级权重合成）}
\label{fig:weight}
\end{figure}

\begin{table}[H]
\centering
\caption{二级指标综合权重（按综合权重降序，仅列出前10项）}
\label{tab:weights}
\begin{tabular}{ccccc}
\toprule
指标 & 类型 & 所属B层 & C层权重 & 综合权重 \\
\midrule
C$_3$ 性价比        & 偏大型 & B$_1$ & 0.2468 & 0.1332 \\
C$_2$ 产品质量      & 偏大型 & B$_1$ & 0.2468 & 0.1332 \\
C$_6$ 店铺信用度    & 偏大型 & B$_1$ & 0.1396 & 0.0753 \\
C$_7$ 售后服务      & 偏大型 & B$_1$ & 0.1396 & 0.0753 \\
C$_8$ 材料费        & 偏小型 & B$_2$ & 0.2499 & 0.0742 \\
C$_{16}$ 销售人员能力 & 偏大型 & B$_3$ & 0.4231 & 0.0692 \\
C$_{10}$ 设备维修费用 & 偏小型 & B$_2$ & 0.1433 & 0.0426 \\
C$_{11}$ 运行费用    & 偏小型 & B$_2$ & 0.1433 & 0.0426 \\
C$_{12}$ 设备折旧费用 & 偏小型 & B$_2$ & 0.1433 & 0.0426 \\
C$_9$ 运输费用      & 偏小型 & B$_2$ & 0.1433 & 0.0426 \\
\bottomrule
\end{tabular}
\end{table}

前三权重为性价比（13.32\%）、产品质量（13.32\%）和店铺信用度（7.53\%），
均属产品综合评价维度，与该维度一级权重占53.96\%的结构吻合。

\subsection{两级模糊合成模型}

\subsubsection{合成算子选择}

模糊综合评价中常用的合成算子有四种\cite{xie2000}：
$M(\wedge,\vee)$（取小取大，主因素决定型）、
$M(\cdot,\vee)$（乘积取大）、
$M(\wedge,\oplus)$（取小有界和）和
$M(\cdot,\oplus)$（加权平均型）。
本文选用\textbf{$M(\cdot,\oplus)$加权平均型}：

\begin{equation}
    b_j = \sum_{i=1}^{n} w_i \cdot r_{ij}, \quad j = 1,\ldots,5
\label{eq:weighted_avg}
\end{equation}

原因：该算子将\textbf{全部}权重与\textbf{全部}隶属度逐项相乘后求和，
不因取小或取大运算而截断信息，适合指标众多、权重分散的综合评价\cite{chen2000}。

\subsubsection{合成流程}

对应于三层指标体系，模糊合成分两级进行：

\textbf{第一级（C$\rightarrow$B）：}
对每个一级指标$B_k$（$k=1,2,3$），从其下属二级指标的隶属度矩阵中提取行构成子矩阵$R_{B_k}$，
与二级权重$W_C^{(k)}$加权合成，得到$1\times 5$的一级评价向量：

\begin{equation}
    B_k = W_{C}^{(k)} \cdot R_{B_k}
\label{eq:first_level}
\end{equation}

\textbf{第二级（B$\rightarrow$A）：}
将三个一级向量组合为总评价矩阵$R_{\text{total}} = [B_1; B_2; B_3]_{3\times 5}$，
再与一级权重$W_B$合成，得到公司整体评价向量$B$（$1\times 5$）：

\begin{equation}
    B = W_{B} \cdot R_{\text{total}}
\label{eq:second_level}
\end{equation}

\textbf{分数转化：}
将$B$与分值向量$S=[100,80,60,40,20]$做内积，将模糊的五维隶属度压缩为单一可比分数：

\begin{equation}
    \text{Score} = B \cdot S^{T} = \sum_{j=1}^{5} b_j \cdot s_j
\label{eq:score}
\end{equation}

\section{实证计算与结果分析}

\subsection{评价总逻辑}

本文的评价策略是\textbf{逐公司独立评价，统一标尺横向比较}：
每家公司均经历"标准化数据 $\rightarrow$ 隶属度矩阵 $\rightarrow$
一级合成 $\rightarrow$ 二级合成 $\rightarrow$ 综合得分"的完整流程，
五级标尺$V$与分值$S$是对所有公司统一适用的计量基准。
以下各节逐步展示公司a的完整计算过程，然后汇总七家公司的结果。

\subsection{隶属度计算}

以公司a的C$_3$性价比指标（$x=87.0$，偏大型）为例，展示边界过渡效应：

\begin{itemize}[itemsep=2pt]
    \item 对$v_1$（优秀），参数$[80,90,100,100]$：87.0在$[80,90)$，
          $\mu = (87.0-80)/(90-80) = 0.700$；
    \item 对$v_2$（良好），参数$[70,80,80,90]$：87.0在$(80,90]$，
          $\mu = (90-87.0)/(90-80) = 0.300$；
    \item 对$v_3$$\sim$$v_5$：均超出覆盖范围，$\mu=0.000$。
\end{itemize}

得到隶属度向量$[0.700, 0.300, 0, 0, 0]$——指标70\%隶属"优秀"，30\%隶属"良好"。
这正是FCE区别于"一刀切"评分的关键：保留了指标值在两个等级边界处的模糊过渡信息。
对全部19个指标重复此操作，得到公司a完整的$19\times 5$隶属度矩阵$R_a$（\cref{tab:membership}）。

\begin{table}[H]
\centering
\caption{公司a的隶属度矩阵$R_a$（19 $\times$ 5）}
\label{tab:membership}
\scriptsize
\begin{tabular}{lccccc}
\toprule
指标 & 优秀 & 良好 & 中等 & 较差 & 很差 \\
\midrule
C$_1$ 产品品牌      & 1.000 & 0.000 & 0.000 & 0.000 & 0.000 \\
C$_2$ 产品质量      & 1.000 & 0.000 & 0.000 & 0.000 & 0.000 \\
C$_3$ 性价比        & 0.700 & 0.300 & 0.000 & 0.000 & 0.000 \\
C$_4$ 产品款式      & 0.910 & 0.090 & 0.000 & 0.000 & 0.000 \\
C$_5$ 产品包装      & 0.030 & 0.970 & 0.000 & 0.000 & 0.000 \\
C$_6$ 店铺信用度    & 1.000 & 0.000 & 0.000 & 0.000 & 0.000 \\
C$_7$ 售后服务      & 0.690 & 0.310 & 0.000 & 0.000 & 0.000 \\
C$_8$ 材料费        & 0.680 & 0.320 & 1.000 & 1.000 & 1.000 \\
C$_9$ 运输费用      & 0.000 & 1.000 & 1.000 & 1.000 & 1.000 \\
C$_{10}$ 设备维修费用 & 0.000 & 1.000 & 1.000 & 1.000 & 1.000 \\
C$_{11}$ 运行费用   & 0.000 & 1.000 & 1.000 & 1.000 & 1.000 \\
C$_{12}$ 设备折旧费用 & 0.070 & 0.930 & 1.000 & 1.000 & 1.000 \\
C$_{13}$ 人员工资   & 0.000 & 0.080 & 0.920 & 1.000 & 1.000 \\
C$_{14}$ 电耗费用   & 0.000 & 1.000 & 1.000 & 1.000 & 1.000 \\
C$_{15}$ 水耗费用   & 0.000 & 1.000 & 1.000 & 1.000 & 1.000 \\
C$_{16}$ 销售人员能力 & 0.000 & 0.800 & 0.200 & 0.000 & 0.000 \\
C$_{17}$ 广告宣传   & 0.000 & 0.320 & 0.680 & 0.000 & 0.000 \\
C$_{18}$ 行业需求   & 0.000 & 0.010 & 0.990 & 0.000 & 0.000 \\
C$_{19}$ 家庭需求   & 0.000 & 0.440 & 0.560 & 0.000 & 0.000 \\
\bottomrule
\end{tabular}
\end{table}

\subsection{一级模糊合成}

以公司a的B$_1$产品综合评价为例，提取C$_1$$\sim$C$_7$的子矩阵$R_{B_1}^a$（$7\times 5$），
与二级权重$W_C^{(1)}$（\cref{tab:b1c}）按\cref{eq:first_level}合成：

\begin{align}
B_1^{a} &= [0.0758, 0.2468, 0.2468, 0.0758, 0.0758, 0.1396, 0.1396] \nonumber \\
        &\quad \cdot \begin{bmatrix}
        1.000&0.000&0.000&0.000&0.000 \\
        1.000&0.000&0.000&0.000&0.000 \\
        0.700&0.300&0.000&0.000&0.000 \\
        0.910&0.090&0.000&0.000&0.000 \\
        0.030&0.970&0.000&0.000&0.000 \\
        1.000&0.000&0.000&0.000&0.000 \\
        0.690&0.310&0.000&0.000&0.000
        \end{bmatrix} \nonumber \\
        &= [0.802, 0.198, 0.000, 0.000, 0.000]
\label{eq:b1}
\end{align}

结果解读：公司a在产品综合评价维度上，80.2\%隶属"优秀"，19.8\%隶属"良好"。
对其余B$_2$、B$_3$和全部公司执行同样计算，一级结果如\cref{tab:first_level}。

\begin{table}[H]
\centering
\caption{七家公司一级模糊评价结果}
\label{tab:first_level}
\small
\begin{tabular}{clccc}
\toprule
公司 & 一级指标 & 评价向量（优秀, 良好, 中等, 较差, 很差） & 最大等级 \\
\midrule
\multirow{3}{*}{a}
    & B$_1$ 产品 & (0.802, 0.198, 0.000, 0.000, 0.000) & 优秀 \\
    & B$_2$ 成本 & (0.046, 0.190, 0.253, 0.255, 0.255) & 较差 \\
    & B$_3$ 销售 & (0.000, 0.480, 0.520, 0.000, 0.000) & 中等 \\
\midrule
\multirow{3}{*}{b}
    & B$_1$ 产品 & (0.314, 0.474, 0.212, 0.000, 0.000) & 良好 \\
    & B$_2$ 成本 & (0.058, 0.182, 0.250, 0.255, 0.255) & 较差 \\
    & B$_3$ 销售 & (0.000, 0.100, 0.663, 0.237, 0.000) & 中等 \\
\midrule
\multirow{3}{*}{c}
    & B$_1$ 产品 & (0.168, 0.562, 0.269, 0.000, 0.000) & 良好 \\
    & B$_2$ 成本 & (0.036, 0.166, 0.253, 0.272, 0.272) & 较差 \\
    & B$_3$ 销售 & (0.499, 0.459, 0.042, 0.000, 0.000) & 优秀 \\
\midrule
\multirow{3}{*}{d}
    & B$_1$ 产品 & (0.787, 0.213, 0.000, 0.000, 0.000) & 优秀 \\
    & B$_2$ 成本 & (0.007, 0.243, 0.250, 0.250, 0.250) & 中等 \\
    & B$_3$ 销售 & (0.478, 0.506, 0.016, 0.000, 0.000) & 良好 \\
\midrule
\multirow{3}{*}{e}
    & B$_1$ 产品 & (0.456, 0.480, 0.064, 0.000, 0.000) & 良好 \\
    & B$_2$ 成本 & (0.101, 0.116, 0.223, 0.268, 0.292) & 很差 \\
    & B$_3$ 销售 & (0.000, 0.041, 0.503, 0.456, 0.000) & 中等 \\
\midrule
\multirow{3}{*}{f}
    & B$_1$ 产品 & (0.113, 0.479, 0.349, 0.060, 0.000) & 良好 \\
    & B$_2$ 成本 & (0.063, 0.093, 0.229, 0.296, 0.319) & 很差 \\
    & B$_3$ 销售 & (0.000, 0.636, 0.364, 0.000, 0.000) & 良好 \\
\midrule
\multirow{3}{*}{g}
    & B$_1$ 产品 & (0.210, 0.315, 0.466, 0.009, 0.000) & 中等 \\
    & B$_2$ 成本 & (0.045, 0.109, 0.230, 0.289, 0.327) & 很差 \\
    & B$_3$ 销售 & (0.000, 0.165, 0.583, 0.227, 0.024) & 中等 \\
\bottomrule
\end{tabular}
\end{table}

七家公司各维度的表现差异通过\cref{fig:radar}的雷达图和多维对比图（\cref{fig:dimension}）可直观呈现。

\begin{figure}[H]
\centering
\begin{subfigure}[b]{0.48\textwidth}
    \centering
    \includegraphics[width=\textwidth]{radar_chart.png}
    \caption{一级指标雷达图}
    \label{fig:radar}
\end{subfigure}
\hfill
\begin{subfigure}[b]{0.48\textwidth}
    \centering
    \includegraphics[width=\textwidth]{dimension_comparison.png}
    \caption{一级指标维度对比}
    \label{fig:dimension}
\end{subfigure}
\caption{七家公司三个一级维度的表现对比}
\label{fig:compare}
\end{figure}

\subsection{二级合成与综合得分}

每家公司的三个B层向量构成$3\times 5$总评价矩阵，与一级权重$W_B$合成：

\begin{equation}
    B = 0.5396 \times B_1 + 0.2970 \times B_2 + 0.1634 \times B_3
\end{equation}

以公司a为例：

\begin{align}
    B_a &= 0.5396 \times [0.802, 0.198, 0.000, 0.000, 0.000] \nonumber \\
        &\quad + 0.2970 \times [0.046, 0.190, 0.253, 0.255, 0.255] \nonumber \\
        &\quad + 0.1634 \times [0.000, 0.480, 0.520, 0.000, 0.000] \nonumber \\
        &= [0.447, 0.242, 0.160, 0.076, 0.076]
\end{align}

解读：公司a整体以44.7\%隶属"优秀"，24.2\%隶属"良好"，16.0\%隶属"中等"，
各7.6\%隶属"较差"和"很差"。

将$B_a$与分值向量$S$做内积，将模糊的多维隶属度转化为单一综合得分：

\begin{align}
    \text{Score}_a &= 0.447 \times 100 + 0.242 \times 80 + 0.160 \times 60
                    + 0.076 \times 40 + 0.076 \times 20 \nonumber \\
                   &= 44.7 + 19.4 + 9.6 + 3.0 + 1.5 = 78.15
\end{align}

七家公司的最终评价结果及排序如\cref{tab:final}，可视化见\cref{fig:score}和\cref{fig:grade}。

\begin{table}[H]
\centering
\caption{七家公司最终模糊评价结果与排序}
\label{tab:final}
\begin{tabular}{cccccc}
\toprule
排名 & 公司 & 最终评价向量（优秀, 良好, 中等, 较差, 很差） & 最大等级 & 得分 \\
\midrule
1 & d & (0.505, 0.270, 0.077, 0.074, 0.074) & 优秀 & 81.14 \\
2 & a & (0.447, 0.242, 0.160, 0.076, 0.076) & 优秀 & 78.15 \\
3 & c & (0.183, 0.428, 0.227, 0.081, 0.081) & 良好 & 71.02 \\
4 & e & (0.276, 0.300, 0.183, 0.154, 0.087) & 良好 & 70.50 \\
5 & b & (0.186, 0.326, 0.297, 0.115, 0.076) & 良好 & 68.66 \\
6 & f & (0.080, 0.390, 0.315, 0.120, 0.095) & 良好 & 64.79 \\
7 & g & (0.127, 0.229, 0.415, 0.128, 0.101) & 中等 & 63.04 \\
\bottomrule
\end{tabular}
\end{table}

\begin{figure}[H]
\centering
\begin{subfigure}[b]{0.48\textwidth}
    \centering
    \includegraphics[width=\textwidth]{score_ranking.png}
    \caption{综合得分排序（含等级分界线）}
    \label{fig:score}
\end{subfigure}
\hfill
\begin{subfigure}[b]{0.48\textwidth}
    \centering
    \includegraphics[width=\textwidth]{grade_distribution.png}
    \caption{评价等级隶属度堆叠分布}
    \label{fig:grade}
\end{subfigure}
\caption{七家公司综合得分排名与等级隶属度分布}
\label{fig:result}
\end{figure}

\subsection{结果分析}

从\cref{tab:final}和\cref{fig:result}可以得出以下结论：

\textbf{第一梯队（75分以上）：}公司d（81.14分）和公司a（78.15分），评价等级均为"优秀"。
公司d的优势在于三个维度的均衡性：产品维度优秀（0.787）、成本维度中等（0.243，七家中最佳）、
销售维度良好（0.506）。结合原始数据（\cref{tab:rawdata}），
d的成本费用指标值（C$_8$$\sim$C$_{15}$）全面低于对手，尤其是材料费（8.6 vs 行业均值$\approx$28）
和水耗费用（0.6），形成了显著的成本护城河。公司a在产品维度与d相当（0.802 vs 0.787），
但成本维度逊于d（B$_2$最大隶属"较差"），因此总得分落后约3分。

\textbf{第二梯队（70分上下）：}公司c（71.02分）和e（70.50分），评价为"良好"。
c凭借销售能力维度"优秀"（0.499）跻身第三，这源于其行业需求（90.1）和家庭需求（86.0）的高分；
但成本维度"较差"制约了总分提升。e的成本维度"很差"（0.292），是七家中成本管控最弱的之一。

\textbf{中等梯队（60$\sim$69分）：}b（68.66）、f（64.79）、g（63.04）。
g是唯一最终等级为"中等"的公司，三大维度均无突出表现。

\textbf{整体规律：}产品综合评价维度（权重53.96\%）几乎决定了公司能否进入第一梯队；
成本费用（29.70\%）是区分名次的关键——低排名公司的共性特征是成本维度评为"很差"；
销售能力（16.34\%）次之但不可忽略，c正是凭此超越e和b夺得第三。

\section{结论与展望}

\subsection{主要结论}

本文以七家公司为研究对象，构建了3一级指标$\times$19二级指标的三层评价体系，
综合运用梯形模糊隶属函数、AHP方根法和$M(\cdot,\oplus)$两级模糊合成模型，
完成了多指标模糊综合评价与排序。主要结论：

\begin{enumerate}[label=(\arabic*),itemsep=2pt]
    \item \textbf{FCE有效处理了等级边界模糊问题。}
          梯形隶属函数将精确指标值转化为五等级隶属度分布，保留了过渡区间信息，
          避免了"非此即彼"刚性分类。如性价比87分时呈70\%优秀+30\%良好的隶属分布，
          精确刻画了边界过渡状态。
    \item \textbf{AHP方法建立了逻辑一致的权重体系。}
          四个判断矩阵的$CR$均小于0.1，通过一致性检验。
          产品综合评价权重最高（53.96\%），性价比与产品质量并列最关键（各13.32\%），
          与企业管理直觉吻合。
    \item \textbf{评价结果合理、可解释。}
          最终排名为d（81.14）$>$ a（78.15）$>$ c（71.02）$>$ e（70.50）$>$
          b（68.66）$>$ f（64.79）$>$ g（63.04）。
          公司d凭借全面均衡的表现获最优评级，各公司的排名与原始数据特征高度一致。
\end{enumerate}

\subsection{不足与展望}

本研究存在以下可改进方向：
（1）隶属函数参数依赖经验设定，未来可引入聚类分析等方法从历史评价数据中自动学习参数；
（2）指标体系当前以产品、成本、销售为主，可扩展纳入财务绩效、创新能力、ESG等维度；
（3）可与TOPSIS、灰色关联分析等方法的排序结果进行一致性检验，进一步验证模型稳健性；
（4）当前实证基于模拟数据，后续可替换为真实企业数据进行检验。

% ============================================================
%  参考文献
% ============================================================
\begin{thebibliography}{99}

\bibitem{wang1983}
汪培庄. 模糊集合论及其应用[M]. 上海: 上海科学技术出版社, 1983.

\bibitem{saaty1980}
Saaty T L. \textit{The Analytic Hierarchy Process: Planning, Priority Setting,
Resource Allocation}[M]. New York: McGraw-Hill, 1980.

\bibitem{xie2000}
谢季坚, 刘承平. 模糊数学方法及其应用[M]. 3版. 武汉: 华中科技大学出版社, 2000.

\bibitem{zadeh1965}
Zadeh L A. Fuzzy sets[J]. \textit{Information and Control}, 1965, 8(3): 338--353.

\bibitem{chen2000}
陈守煜. 工程模糊集理论与应用[M]. 北京: 国防工业出版社, 2000.

\bibitem{saaty2008}
Saaty T L. Decision making with the analytic hierarchy process[J].
\textit{International Journal of Services Sciences}, 2008, 1(1): 83--98.

\bibitem{hwang1981}
Hwang C L, Yoon K. \textit{Multiple Attribute Decision Making: Methods and
Applications}[M]. Berlin: Springer-Verlag, 1981.

\end{thebibliography}

\end{document}
